% List of abbreviations. Use \ac{KEY} in the text: the first use prints
% "long form (SHORT)", later uses print "SHORT". \acp{KEY} = plural.
% Only the abbreviations you actually use are printed (option printonlyused).
% The optional argument of the environment is the longest short form (sets the column width).

\chapter*{List of abbreviations}
\addcontentsline{toc}{chapter}{List of abbreviations}

\begin{acronym}[FB-eTRCA]
  \acro{ALS}{amyotrophic lateral sclerosis}
  \acro{BCI}{brain--computer interface}
  \acro{CCA}{canonical correlation analysis}
  \acro{CI}{continuous integration}
  \acro{EEG}{electroencephalography}
  \acro{eTRCA}{ensemble task-related component analysis}
  \acro{FBCCA}{filter bank canonical correlation analysis}
  \acro{FBeTRCA}[FB-eTRCA]{filter bank ensemble task-related component analysis}
  \acro{GUI}{graphical user interface}
  \acro{GPU}{graphics processing unit}
  \acro{ITR}{information transfer rate}
  \acro{JSON}{JavaScript Object Notation}
  \acro{LDA}{linear discriminant analysis}
  \acro{LM}{language model}
  \acro{LSL}{Lab Streaming Layer}
  \acro{NIRS}{near-infrared spectroscopy}
  \acro{NLP}{natural language processing}
  \acro{PSD}{power spectral density}
  \acro{SNR}{signal-to-noise ratio}
  \acro{SSVEP}{steady-state visual evoked potential}
  \acro{TRCA}{task-related component analysis}
\end{acronym}
